\chapter{Resumen}

El objetivo de la tesis es caracterizar tres aspectos de los asentamientos humanos en México: su extensión geográfica, su distribución de tamaño y la estructura de sus redes de calles. Para ello desarrollo y aplico métodos matemáticos, específicamente de geometría computacional, teoría de las gráficas y las redes, probabilidad y estadística. En el primer capítulo propongo un algoritmo de delimitación de asentamientos humanos basado en la gráfica de intersección del casco convexo de polígonos que representan áreas construidas. Muestro que los límites generados por el algoritmo reflejan mejor la extensión geográfica de las ciudades que sus delimitaciones administrativas o económicas. En el segundo capítulo elaboro un estado del arte de los métodos estadísticos para verificar si una muestra empírica sigue una distribución de ley de potencia. Aplico estos métodos para caracterizar la distribución de tamaño de ciudades de México, Estados Unidos, Canadá, Chile, Colombia y Uruguay. Encuentro que la distribución del tamaño de las ciudades en México ha seguido una ley de potencia durante los últimos 20 años, la cual se extiende hasta concentraciones urbanas pequeñas. También identifico que la distribución de tamaño de los tres países de Norteamérica sigue una ley de potencia, tanto individualmente como en conjunto. En el tercer capítulo abordo teórica y computacionalmente algunas propiedades de los $\beta$-esqueletos. Demuestro que las teselaciones de polígonos cíclicos son $\beta$-esqueletos. Además, observo que los $\beta$-esqueletos son capaces de aproximar bien teselaciones regulares con vértices perturbados moderadamente. Aplico dos teoremas de la literatura para calcular el grado y la longitud promedio de los $\beta$-esqueletos de puntos aleatorios. En el último capítulo analizo 11,091 redes de calles mediante índices \textit{topológicos} y geométricos propuestos en la literatura, así como con el grado de aproximación de los $\beta$-esqueletos. Una conclusión obtenida es que el nivel de aproximación de los $\beta$-esqueletos puede ser utilizado como una alternativa para evaluar la similitud de una red de calles a una malla rectangular.